%!TeX root = ../incremental_SS_Translation.tex

\chapter{Śārīrasthāna 8:  Phlebotomy}


\section{Introduction}
\texttt{Draft text by Jan Gerris.
    The Sanskrit text of this adhyāya is a lightly 
    edited transcript of NAK 5-333,
    waiting for full critical attention.}

\section{Literature} 

Meulenbeld offered an annotated overview of this chapter and a
bibliography of earlier scholarship to 2002 and, in his notes, citations
of the parallel passages in the \CS.\fvolcite{IA}[257--258]{meul-hist}   

The first draft of the following translation was prepared by Jan
Gerris.
% and later revised in collaboration Philipp Maas, John Nemec,
%Dominik Wujastyk, Madhusudan Rimal, Kenneth Zysk and Andrey Klebanov.

\newpage

\section{Translation}

\begin{translation}
    \bigskip       

\begin{tt}
    \raggedright
    \bigskip
    \marginpar{Draft tr.\ from here}

\item[1]
Now, therefore, we shall expound the chapter on the body regarding the procedure for venesection.

\item[2]   vulgate only

\item[3]
In this regard, one should not pierce the veins of children, the elderly, the timid, those who are dehydrated, those emaciated by chest wounds, those weakened by alcohol, travel, or sexual indulgence, those who have not undergone enema treatment, those who have undergone emesis or purgation, those who have received dental enemas, the impotent, pregnant women, and those afflicted by cough, asthma, consumption, high fever, diarrhoea, convulsions, paralysis, or fasting.
What is the reason for this? Indeed, from piercing these individuals, there occurs confusion of the senses or wasting.
In those with bleeding disorders, due to the excessive loss of blood, there is provocation of Vata or death.
For pregnant women, however, there is wasting of the foetus, miscarriage, or the fetus may have deficient limbs or a short lifespan.

\item[3a]
And there is a verse on this subject:
By piercing the veins of those for whom bloodletting is prohibited, an increase of the disease or death can occur even for the remaining living beings.

\item[3b]
And those [veins] which are seen or unseen, and those which are not controlled or not raised, are considered unfit for piercing.

\item[4]
The disorders curable by bloodletting have been mentioned previously; in those cases, when they are not yet suppurated, and in others not yet experienced, one should pierce the vein according to repeated practice and without causing an emergency.

\item[5]
Even for those generally prohibited, venesection is not forbidden in cases of poisoning or acute infection.

\item[6]
In that case, having prepared the patient who has been anointed and fomented, who has eaten food that is mostly liquid and contrary to the humours or has drunk gruel, and having placed him at the proper time in a sitting or standing position without causing distress to his vital breaths.
Having secured the body part with any one of cloth, linen, leather, inner bark, or creepers, neither too tightly nor too loosely, and taking the instrument according to the rules, one should pierce the prescribed vein.

\item[7]
Regarding the vein to be pierced: Venesection should never be performed when it is too cold, too hot, too windy, or on a cloudy day, nor should it ever be done to a healthy person.

\item[8]
In that case, having seated the man whose vein is to be pierced facing the sun on a seat one Aratni high, having placed his elbows upon his two flexed knee-joints, and having placed his hands—with fists made by hiding the thumbs inside—upon the other [Manya muscles], and having thrown the binding cloth over the neck, the physician should have another person standing behind hold the two ends of the cloth with an upturned left hand, and for the purpose of raising the vein and draining the blood, he should apply pressure to the apparatus according to the proper measure.
The patient should fill his mouth with air.
This is the method for raising the veins of the head without mechanical restraint.
Having placed the patient with his thighs steady and his foot slightly curved on a high place, and having properly tied the ligature above the area to be pierced, one should pierce the vein of the foot.
For a patient seated comfortably with his hand placed on a proper seat and his thumb hidden in his fist, one should pierce the vein of the hand.
In cases of sciatica and brachial neuralgia, the knee and elbow should be flexed; for veins in the hips, back, and shoulders, the back should be raised; for veins in the abdomen and chest, the chest should be expanded and the body stretched out; and for veins in the sides, the patient should be supported by his two arms.
For a vein in the penis, the penis should be bent down.
For a vein under the tongue, the tongue should be raised and its tip bitten.
When opening a vein in the palate or in the roots of the teeth, the mouth should be kept wide open.
Similarly, one should intelligently devise other mechanical means and methods for raising the veins, considering the patient's physical condition and the nature of the disease.

\item[9]
In muscular areas, one should apply the Vrihimukha instrument to the depth of a barley-corn or a grain of rice.
Otherwise, over the bones, one should puncture with the Kutharika to the depth of half a barley-corn.

\item[10]
And there are verses on this subject:
In the rainy season, one should perform venesection when the sky is clear of clouds, and during the summer, in the cool part of the day; at midday during the winter—these are remembered as the three times for using the instrument.


\item[11]
One should recognize a vein as well-pierced if the blood flows with a stream like that of twice-churned buttermilk, and if it stops of its own accord after being held for a Muhūrta.

\item[12]
Just as yellow fluid first flows from Kusumbha flowers, so too the vitiated blood flows out first when the veins are pierced.

\item[13]  vulgate only

\item[14]
In a patient who is weakened, has many accumulated doṣas, or is afflicted by fainting, the blood should be drained again in the afternoon, or on the following day if the Vāyu is provoked.

\item[15]
A wise physician should leave some of the blood containing residual doṣas; he should not cause excessive bleeding, but should overcome the remainder with soothing remedies.

\item[16]  vulgate only

\item[17]
In diseases such as burning feet, numbness of feet, whitlow, Vātarakta, Vātakaṇṭaka, Vipādikā, and cracked heels, one should pierce the vein two finger-breadths above the Kṣipra-marma with a Vrīhimukha instrument.
In cases of elephantiasis, as previously stated; in Kroṣṭukaśīrṣa and Vāta-pains, the vein located in the thigh; in Apacī, four finger-breadths above the ankle joint; two finger-breadths below the Indrabasti; in sciatica, four finger-breadths above the knee joint; in goitre, the vein situated at the hip-joint; in Pravāhikā with pain, around the pelvis; and in Parivartikā, Upadaṃśa, and Śūka-disorders, in the middle of the penis.
Specifically, in cases of splenic disorders, this same vein should be opened in the middle of the left arm, on the inner side of the elbow joint; in cases of Yakṛddālya, the same vein in the right arm.
Authorities also prescribe this same vein for cough and asthma. 
In cases of sciatica or Vipañcī, the vein situated between the left side, the armpit, and the breast should be opened; likewise in cases of abscess and pain in the sides.
This same vein is used in Kaphodara.
Some say that in cases of atrophy of the arm and Avabāhuka, the vein in the middle of the arm on the right side should be opened; in tertian fever, the vein in the middle of the shoulders or the middle of the Trika joint.
In quartan fever, the veins located at the shoulder joint or on either side should be opened.
In cases of epilepsy, the vein situated in the middle of the jaw joint should be opened.
In cases of insanity, the veins situated at the junction of the temples.
In diseases of the tongue, the veins on the under-surface of the tongue.
In dental diseases, at the root of the teeth.
In the palate and in diseases of the palate.
In ear-ache and other diseases of the ears, the veins all around the upper part of the ears.
In diseases of the nose, at the tip of the nose for the perception of smells.
In eye diseases such as Timira and Akṣipāka, the veins near the nose, on the forehead, or at the outer corner of the eye.
These same veins should also be opened in diseases such as headaches and Adhimantha.

\item[18]
In this context, there are twenty types of defective venesections: durviddha (badly pierced), atividdha (excessivelypierced), kuñcita (contracted), piccita (crushed),  kuṭṭita (lacerated), aprasrutā (non-flowing), anudīrṇā (insufficiently opened), anteviddha (pierced at the edge),  pariśuṣkā (dried up), kūṇitā (narrowed), vepitā (trembling), 
Anutthitaviddha (pierced without being raised), avyadhyaviddha (pierced where it should not be), śastrahata (wounded by the instrument), tiryagviddha (obliquely pierced), aviddha (improperly pierced), vidrutā (unsteady), dhenukā (repeatedly bleeding), punar-punarviddha (pierced again and again), and those pierced in the vital spots, veins, ligaments, bones, and joints.

\item[19]
Among these, when blood does not flow clearly due to the application of a very fine instrument, and pain and swelling occur, that is called durviddha.
When the piercing is beyond the proper measure, blood either enters internally or there is excessive discharge; that is called atividdha, and the same occurs in kuñcita.
A vein that has become flattened and broad due to being crushed by a blunt instrument is called piccita; kuṭṭita is that which is struck many times without reaching the interior.
That which does not yield blood due to cold, fear, or fainting is called aprasrutā; that which is pierced with a sharp, large-mouthed instrument but is not sufficiently opened is called anudīrṇā; and that which yields very little blood is called anteviddha.
A vein in a patient with depleted blood that is filled with wind and dried up is called pariśuṣkā.
A vein that is only a quarter-part incised and from which blood slightly issues is called kūṇitā.
When there is confusion of blood flow in a vein that trembles due to being bandaged in a wrong place, it is called vepitā.
The same result occurs when a vein is not properly raised; those that should not be operated upon with an instrument are called avyadhya.
A vein that, after being cut, has excessive blood flow and hinders the procedure is called śastrahata.
A vein into which the instrument is applied slantingly, leaving it partially cut, is called tiryagviddha.
A vein that is wounded several times due to the improper application of the instrument is called aviddha.
A vein that is pierced while it is unsteady is called vidruta.|
A vein that bleeds excessively due to the area being frequently pressed and the incision not healing is called dhenukā.
Piercing again and again, or piercing in the veins, vital spots, ligaments, bones, and joints, results in pain, swelling, deformity, or death.



\item[20]
And here are the verses on this subject: There is no one perfectly skilled in veins, for they are unsteady by nature; they move about like fish, therefore one should strike them with great care.

\item[21]
When a surgical instrument is applied to the body by an ignorant person...
These calamities and many other complications occur.

\item[22]
Diseases are quickly pacified by proper venesection, in a way that is not possible through the application of oleation and other procedures, nor even by medicinal plasters.

\item[23]
Venesection is prescribed for the sake of treatment in the field of surgery, just as a properly administered enema is in internal medicine.

\item[24]
In that regard, men who have undergone oleation, sweating, vomiting, purging, oily enemas, decoction enemas, or venesection must avoid certain things.
Some consider that one should avoid anger, exhaustion, sexual intercourse, sleeping during the day, speaking loudly, physical exercise, jolting from vehicles, excessive sitting, standing, or walking, cold water, wind, sun, incompatible foods, unwholesome food, and indigestion until one regains strength.
We shall describe the details of these matters later on.

\item[25]
And there is a verse on this subject: 
Vitiated blood is extracted by means of venesection, animal horns, gourds, leeches, or scratching, following the previously mentioned methods according to the depth.


Thus ends the eighth chapter in the section on the body.
\end{tt}
\end{translation}